\chapter[Fluid Mechanics]{Fluid Mechanics: Balance Laws, Navier-Stokes, And Stability}
\label{chap:fluid}

\section{What Is A Fluid?}
\label{sec:what-is-fluid}

Fluid mechanics is continuum mechanics with a special constitutive idea: the
material does not resist a static shear deformation. A solid can sit in a
sheared shape and store elastic energy; a simple fluid responds to shear through
stress proportional to the rate of deformation. This chapter therefore treats
fluids on their own terms. Elastic and deforming bodies are separated into the
next chapter.

The basic fields are functions on physical space:
\[
  \rho(x,t),\qquad u(x,t),\qquad p(x,t),
\]
where \(\rho\) is mass density, \(u\) is velocity, and \(p\) is pressure. A
fluid particle can still be labelled by a material coordinate \(X\), with
\[
  x=\varphi(X,t),\qquad u(\varphi(X,t),t)=\partial_t\varphi(X,t).
\]
The Eulerian view keeps \(x\) fixed and watches fluid pass through; the
Lagrangian view follows the particle labelled by \(X\).

\begin{table}[htbp]
\centering
\caption{Fluid-mechanics notation. The symbol \(p\) denotes pressure in this
chapter, not canonical momentum.}
\label{tab:fluid-notation}
\small
\begin{tabular}{p{0.18\textwidth}p{0.48\textwidth}p{0.25\textwidth}}
\toprule
Symbol & Meaning & Type\\
\midrule
\(X\) & material label of a fluid particle & coordinate in reference region\\
\(\varphi(X,t)\) & flow map & material-to-spatial map\\
\(F^i{}_A\) & \(\partial\varphi^i/\partial X^A\) & deformation gradient of flow map\\
\(J\) & \(\det F\) & local volume ratio\\
\(x\) & Eulerian spatial position & point in physical space\\
\(\rho(x,t)\) & mass density & scalar field\\
\(u(x,t)\) & velocity & vector field\\
\(p(x,t)\) & pressure & scalar field\\
\(V(t)\) & material volume & moving region in space\\
\(n\) & outward unit normal to a surface & vector field on boundary\\
\(b(x,t)\) & body force per unit mass & vector field\\
\(t(n)\) & surface traction on normal \(n\) & vector field\\
\(\sigma\) & Cauchy stress tensor & \(3\times3\) tensor field\\
\(D u/Dt\) & material acceleration & \(\partial_tu+u\cdot\nabla u\)\\
\(D\) & rate-of-strain tensor & \(\frac12(\nabla u+\nabla u^T)\)\\
\(W\) & antisymmetric spin tensor & \(\frac12(\nabla u-\nabla u^T)\)\\
\(\omega\) & vorticity & \(\nabla\times u\)\\
\(\mu\) & dynamic viscosity & \(ML^{-1}T^{-1}\)\\
\(\nu\) & kinematic viscosity & \(\mu/\rho\)\\
\(\lambda\) & second viscosity coefficient & stress coefficient\\
\(U,L\) & characteristic velocity and length & nondimensionalization scales\\
\(\mathrm{Re}\) & Reynolds number & \(UL/\nu\)\\
\(\Omega\) & fixed spatial domain in energy estimates & subset of \(\R^3\)\\
\(\gamma(t)\) & closed material loop & curve transported by the flow\\
\(h,\Phi\) & specific enthalpy and conservative body-force potential & Kelvin theorem\\
\(V(y),G\) & laminar profile and pressure-gradient magnitude & shear-flow sections\\
\(\alpha,\sigma,\phi(y)\) & wave number, growth rate, normal-mode amplitude & Orr-Sommerfeld section\\
\(\Gamma_i\) & point-vortex circulation & vortex strength\\
\bottomrule
\end{tabular}
\end{table}

\begin{assumption}[Continuum fluid]
The molecular scale is much smaller than the length scales being modelled, so
density, velocity, pressure, and stress are treated as smooth fields. Across
shocks, contact discontinuities, or interfaces, the differential equations must
be replaced by weak forms and jump conditions.
\end{assumption}

\paragraph{Roadmap.}
The derivation is deliberately staged. First we learn how to differentiate
quantities attached to moving fluid parcels. Then we convert integral balance
laws into local partial differential equations for mass and momentum. Only
after those kinematic and balance statements are in place do we introduce the
Newtonian constitutive law that turns Cauchy's momentum equation into the
Navier-Stokes equation. The later sections use that equation in three ways:
exact laminar solutions, stability analysis, and Hamiltonian point-vortex
reduction.

The staging matters because the Navier-Stokes equation is not a single
assumption. Its inertial term is kinematics, its pressure and stress terms come
from momentum balance, and its viscous term comes from the Newtonian
constitutive law. Separating these ingredients makes it clear which part of the
model changes when one studies inviscid flow, non-Newtonian fluids, elasticity,
or low-Reynolds-number swimming.

\paragraph{Derivation audit trail.}
Every continuum equation below should be read with its conservation statement
attached. The material derivative comes from following one labelled parcel.
The Reynolds transport theorem converts derivatives of moving integrals into
Eulerian fields. Mass conservation comes from conserving mass in every material
volume. Cauchy's equation comes from momentum balance and the existence of a
traction stress tensor. Navier-Stokes appears only after the Newtonian stress
law is imposed. This audit trail is the best protection against treating
\(\rho Du/Dt=-\nabla p+\mu\Delta u\) as an isolated formula rather than the
endpoint of several modelling choices.

\section{Material Derivative}
\label{sec:material-derivative}

If \(f(x,t)\) is a scalar or vector field, its value along a particle path
\(x(t)=\varphi(X,t)\) changes as
\[
  \frac{\dd}{\dd t}f(x(t),t)
  =
  \partial_t f(x(t),t)+\dot x(t)\cdot\nabla f(x(t),t).
\]
Since \(\dot x=u(x(t),t)\), the material derivative is
\[
  \boxed{
  \frac{D f}{D t}
  =
  \partial_t f+u\cdot\nabla f.
  }
\]
For the velocity field itself,
\[
  \frac{D u}{D t}
  =
  \partial_tu+(u\cdot\nabla)u
\]
is the particle acceleration. The nonlinear term is kinematic: it appears even
before any force law is chosen, because the velocity field is sampled along a
moving trajectory.

\section{Reynolds Transport Theorem}
\label{sec:reynolds-transport}

Let \(V(t)\) be a material volume: it moves with the fluid, so its boundary
velocity is \(u\). For a scalar field \(f\), we want the time derivative of
\[
  I(t)=\int_{V(t)} f(x,t)\,\dd^3x.
\]
There are two sources of change. The value of \(f\) can change along each
particle path, and the volume element occupied by neighboring particles can
expand or contract. The Reynolds transport theorem is exactly the formula that
keeps both effects.

\begin{figure}[htbp]
\centering
\begin{tikzpicture}[scale=1.0, >=Latex]
  \draw[->] (-0.4,0) -- (6.8,0) node[right] {\(x^1\)};
  \draw[->] (0,-0.4) -- (0,4.0) node[above] {\(x^2\)};
  \draw[blue!70!black, thick, fill=blue!5]
    (1.0,0.9) .. controls (1.8,0.35) and (3.2,0.55) .. (3.85,1.15)
    .. controls (4.6,1.85) and (4.15,3.0) .. (3.0,3.2)
    .. controls (1.75,3.45) and (0.75,2.45) .. (1.0,0.9);
  \draw[red!75!black, thick, ->] (1.3,1.0) -- (1.75,0.55);
  \draw[red!75!black, thick, ->] (3.9,1.45) -- (4.5,1.55);
  \draw[red!75!black, thick, ->] (2.95,3.15) -- (3.15,3.75);
  \draw[red!75!black, thick, ->] (0.9,2.4) -- (0.45,2.75);
  \node[blue!70!black] at (2.45,2.0) {\(V(t)\)};
  \node[red!75!black] at (5.25,2.85) {boundary velocity \(u\)};
  \draw[red!75!black, ->] (4.75,2.65) -- (3.25,3.45);
  \node[align=center] at (5.15,0.75)
    {material points stay\\inside the moving volume};
\end{tikzpicture}
\caption{A material volume is carried by the fluid velocity. Its boundary
velocity is \(u\), so the time derivative of an integral over \(V(t)\) must
include both the material change of the integrand and the volume change
measured by \(\nabla\cdot u\).}
\label{fig:reynolds-material-volume}
\end{figure}

\begin{derivation}[Transport through a moving volume]
Write \(V(t)=\varphi(U,t)\) for a fixed material region \(U\). Let
\[
  F^i{}_A=\frac{\partial\varphi^i}{\partial X^A},
  \qquad
  J=\det F.
\]
Changing variables gives
\[
  I(t)=\int_U f(\varphi(X,t),t)J(X,t)\,\dd^3X.
\]
Differentiate:
\[
  \dot I
  =
  \int_U
  \left(
    \frac{D f}{D t}J+f\dot J
  \right)\dd^3X.
\]
The Jacobian identity is
\[
  \dot J=J\nabla\cdot u.
\]
One way to see it is Jacobi's formula:
\[
  \dot J=J\,\tr(F^{-1}\dot F),
\]
and \(\dot F^i{}_A=\partial_Au^i(\varphi,t)\), so
\(\tr(F^{-1}\dot F)=\partial_i u^i\). Returning to spatial variables,
\[
  \boxed{
  \frac{\dd}{\dd t}\int_{V(t)} f\,\dd^3x
  =
  \int_{V(t)}
  \left(
    \frac{D f}{D t}+f\nabla\cdot u
  \right)\dd^3x
  =
  \int_{V(t)}
  \left(
    \partial_t f+\nabla\cdot(fu)
  \right)\dd^3x.
  }
\]
\end{derivation}

The theorem separates two effects: the field \(f\) changes along particles, and
the material volume expands or contracts.

\section{Conservation Of Mass}
\label{sec:mass-conservation}

Mass in a material volume is constant:
\[
  \frac{\dd}{\dd t}\int_{V(t)}\rho\,\dd^3x=0.
\]
Using the transport theorem,
\[
  \int_{V(t)}
  \left(
    \partial_t\rho+\nabla\cdot(\rho u)
  \right)\dd^3x=0.
\]
Because this holds for every material volume,
\[
  \boxed{
  \partial_t\rho+\nabla\cdot(\rho u)=0.
  }
\]
Equivalently,
\[
  \frac{D\rho}{Dt}+\rho\nabla\cdot u=0.
\]
For constant-density incompressible flow,
\[
  \boxed{\nabla\cdot u=0.}
\]

\section{Cauchy Momentum Balance}
\label{sec:cauchy-momentum}

Newton's second law for a material volume \(V(t)\) is
\[
  \frac{\dd}{\dd t}\int_{V(t)}\rho u\,\dd^3x
  =
  \int_{V(t)}\rho b\,\dd^3x
  +
  \int_{\partial V(t)} t(n)\,\dd S,
\]
where \(b\) is body force per unit mass and \(t(n)\) is surface traction on a
surface whose outward unit normal is \(n\).

\begin{assumption}[Cauchy stress hypothesis]
The traction is local and linear in the normal:
\[
  t(n)=\sigma n,
\]
where \(\sigma_{ij}(x,t)\) is the Cauchy stress tensor. Balance of angular
momentum, in the absence of couple stresses, implies
\[
  \sigma=\sigma^T.
\]
\end{assumption}

\begin{derivation}[Local momentum equation]
Apply the transport theorem to \(f=\rho u\). Using mass conservation, the
left-hand side becomes
\[
\begin{aligned}
  \frac{\dd}{\dd t}\int_{V(t)}\rho u\,\dd^3x
  &=
  \int_{V(t)}
  \left(
    \partial_t(\rho u)+\nabla\cdot(\rho u\otimes u)
  \right)\dd^3x\\
  &=
  \int_{V(t)}
  \rho\left(\partial_tu+u\cdot\nabla u\right)\dd^3x.
\end{aligned}
\]
The traction term is converted by the divergence theorem:
\[
  \int_{\partial V(t)}\sigma n\,\dd S
  =
  \int_{V(t)}\nabla\cdot\sigma\,\dd^3x.
\]
Therefore
\[
  \int_{V(t)}
  \left[
    \rho\frac{D u}{D t}
    -
    \nabla\cdot\sigma
    -
    \rho b
  \right]\dd^3x=0.
\]
Since \(V(t)\) is arbitrary,
\[
  \boxed{
  \rho\frac{D u}{D t}=\nabla\cdot\sigma+\rho b.
  }
\]
\end{derivation}

This equation is not yet a closed fluid equation. It is a balance law. Closure
comes from a constitutive relation for \(\sigma\).

\section{Newtonian Stress And The Navier-Stokes Equation}
\label{sec:navier-stokes-derivation}

The balance laws so far are universal for continua, but they are not yet a
closed fluid model: the stress tensor \(\sigma\) is still unknown. A
constitutive law supplies the missing relation between stress and motion. For a
simple Newtonian fluid, the stress may depend on the instantaneous rate of
deformation but not on the accumulated shear history. This is the point at
which fluids separate from elastic solids.

The stress of a simple Newtonian fluid is assumed to depend locally and
linearly on the instantaneous rate of deformation, not on the deformation
itself. The velocity gradient splits into symmetric and antisymmetric parts:
\[
  \nabla u=D+W,
  \qquad
  D_{ij}=\frac12(\partial_i u_j+\partial_j u_i),
  \qquad
  W_{ij}=\frac12(\partial_i u_j-\partial_j u_i).
\]
The antisymmetric part \(W\) represents local rigid rotation. It should not
produce viscous stress in an isotropic simple fluid. This statement is a
constitutive assumption with three ingredients:
\[
\begin{array}{ll}
\text{locality:} & \sigma(x,t) \text{ depends on } u \text{ through }
  \nabla u(x,t),\\
\text{objectivity:} & \text{superposed rigid rotations do not create stress},\\
\text{isotropy:} & \text{there is no preferred material direction}.
\end{array}
\]
Objectivity removes \(W\). Isotropy says that a linear map from the symmetric
tensor \(D\) to a symmetric stress can only use the deviatoric part of \(D\)
and its trace. The most general isotropic linear relation is therefore
\[
  \boxed{
  \sigma=-pI+2\mu D+\lambda(\nabla\cdot u)I.
  }
\]
Here \(p\) is pressure, \(\mu\geq0\) is shear viscosity, and \(\lambda\) is the
second viscosity coefficient. For constant \(\mu,\lambda\),
\[
\begin{aligned}
  (\nabla\cdot\sigma)_i
  &=
  -\partial_i p
  +
  \mu\partial_j(\partial_i u_j+\partial_j u_i)
  +
  \lambda\partial_i(\partial_k u_k)\\
  &=
  -\partial_i p
  +
  \mu\Delta u_i
  +
  (\lambda+\mu)\partial_i(\nabla\cdot u).
\end{aligned}
\]
Thus the compressible Navier-Stokes momentum equation is
\[
  \boxed{
  \rho
  \left(
    \partial_tu+u\cdot\nabla u
  \right)
  =
  -\nabla p+\mu\Delta u+(\lambda+\mu)\nabla(\nabla\cdot u)+\rho b.
  }
\]
Together with mass conservation, an equation of state, and an energy or
temperature equation when needed, this describes a Newtonian compressible fluid.

For incompressible flow with constant \(\rho\),
\[
  \nabla\cdot u=0,
\]
and the pressure is best understood as the Lagrange multiplier enforcing this
constraint. The equations reduce to
\[
  \boxed{
  \begin{aligned}
  \rho(\partial_tu+u\cdot\nabla u)&=-\nabla p+\mu\Delta u+\rho b,\\
  \nabla\cdot u&=0.
  \end{aligned}
  }
\]

Read the incompressible equations term by term. The material acceleration
\(\partial_tu+u\cdot\nabla u\) is the acceleration of a moving fluid particle.
The pressure gradient enforces the volume-preserving constraint and transmits
normal stress. The Laplacian \(\nu\Delta u\) diffuses momentum. The body force
adds external acceleration. This interpretation prevents a common mistake:
pressure is not specified independently in incompressible flow; it is solved
simultaneously with \(u\) so that \(\nabla\cdot u=0\) remains true.
Taking the divergence of the incompressible momentum equation makes this
constraint role explicit:
\[
  -\Delta p
  =
  \rho\,\nabla\cdot(u\cdot\nabla u)-\rho\,\nabla\cdot b,
\]
with boundary conditions determined by the normal component of the momentum
equation at the wall. The pressure is therefore a nonlocal field: it adjusts
instantaneously, within the incompressible model, to project the acceleration
back onto divergence-free velocity fields.

\section{Nondimensionalization And Reynolds Number}
\label{sec:reynolds-number}

Let \(U\) be a typical speed, \(L\) a typical length, and define
\[
  x=Lx',\qquad
  t=\frac{L}{U}t',\qquad
  u=Uu',\qquad
  p=\rho U^2p'.
\]
Dropping primes, the incompressible Navier-Stokes equation without body force
becomes
\[
  \partial_tu+u\cdot\nabla u
  =
  -\nabla p+\frac1{\mathrm{Re}}\Delta u,
  \qquad
  \nabla\cdot u=0,
\]
where
\[
  \boxed{\mathrm{Re}=\frac{UL}{\nu}},\qquad \nu=\frac{\mu}{\rho}.
\]
The Reynolds number compares inertial transport to viscous diffusion. Small
\(\mathrm{Re}\) flow is viscously dominated and often reversible in time after
the velocity is reversed. Large \(\mathrm{Re}\) flow can carry disturbances
long distances before viscosity damps them.

\section{Energy Identity And Viscous Dissipation}
\label{sec:fluid-energy-identity}

For a smooth incompressible solution in a fixed domain \(\Omega\), assume either
periodic boundary conditions or no-slip walls \(u=0\) on \(\partial\Omega\).
Dot the Navier-Stokes equation with \(u\) and integrate:
\[
  \rho\int_\Omega u\cdot\partial_tu\,\dd^3x
  +
  \rho\int_\Omega u\cdot(u\cdot\nabla u)\,\dd^3x
  =
  -\int_\Omega u\cdot\nabla p\,\dd^3x
  +
  \mu\int_\Omega u\cdot\Delta u\,\dd^3x.
\]
The nonlinear term is a divergence:
\[
  u\cdot(u\cdot\nabla u)=u\cdot\nabla\left(\frac12\norm u^2\right)
  =
  \nabla\cdot\left(\frac12\norm u^2u\right),
\]
because \(\nabla\cdot u=0\). The pressure term is also a divergence:
\[
  u\cdot\nabla p=\nabla\cdot(pu).
\]
The boundary conditions remove both boundary fluxes. Integration by parts gives
\[
  \int_\Omega u\cdot\Delta u\,\dd^3x
  =
  -\int_\Omega |\nabla u|^2\,\dd^3x.
\]
Hence
\[
  \boxed{
  \frac{\dd}{\dd t}
  \int_\Omega \frac12\rho\norm u^2\,\dd^3x
  =
  -\mu\int_\Omega |\nabla u|^2\,\dd^3x\leq0.
  }
\]
Viscosity irreversibly converts organized kinetic energy into heat. This
identity is also the first stability estimate for laminar flows: disturbances
must draw energy from mean shear fast enough to beat viscous dissipation.

\section{Vorticity And Circulation}
\label{sec:vorticity-circulation}

The vorticity is
\[
  \omega=\nabla\times u.
\]
Taking the curl of the incompressible Navier-Stokes equation gives
\[
  \boxed{
  \partial_t\omega+u\cdot\nabla\omega
  =
  \omega\cdot\nabla u+\nu\Delta\omega.
  }
\]
The term \(\omega\cdot\nabla u\) stretches and tilts vortex lines in three
dimensions. In two dimensions, vorticity is perpendicular to the plane and the
stretching term vanishes:
\[
  \partial_t\omega+u\cdot\nabla\omega=\nu\Delta\omega.
\]
This contrast is central. Two-dimensional inviscid vorticity is transported
like a scalar dye, which creates many conserved quantities. Three-dimensional
vorticity can be stretched by the flow itself, allowing vortex intensification
and making the analysis of Navier-Stokes much harder.
For \(\nu=0\), vorticity is transported by the flow.

\begin{proposition}[Kelvin circulation theorem]
Let \(\gamma(t)\) be a closed material loop moving with an inviscid barotropic
fluid under conservative body forces. Then
\[
  \frac{\dd}{\dd t}\oint_{\gamma(t)}u\cdot \dd x=0.
\]
\end{proposition}

\begin{proof}
Parametrize the loop by \(s\mapsto x(s,t)\), with
\[
  \partial_t x(s,t)=u(x(s,t),t).
\]
Then
\[
  \frac{\dd}{\dd t}\oint u\cdot \partial_sx\,\dd s
  =
  \oint
  \left(
    \frac{D u}{D t}\cdot \partial_sx
    +
    u\cdot\partial_su
  \right)\dd s.
\]
For inviscid barotropic flow with conservative potential \(\Phi\),
\[
  \frac{D u}{D t}=-\nabla h-\nabla\Phi,
\]
where \(h\) is specific enthalpy. The integrand becomes
\[
  \partial_s\left(
    -h-\Phi+\frac12\norm u^2
  \right).
\]
The integral of a single-valued total derivative around a closed loop is zero.
\end{proof}

\section{Laminar Flow Solutions}
\label{sec:laminar-solutions}

Laminar solutions are exact, organized solutions of Navier-Stokes. They are
important not because all flows are laminar, but because stability theory asks
which organized solutions persist under disturbance.
Each solution follows from three choices: geometry, boundary conditions, and
pressure gradient or wall motion. Changing any one of these choices changes the
profile and the instability problem.

\subsection{Plane Couette Flow}
\label{subsec:plane-couette}

Two infinite plates are separated by distance \(h\). The lower plate is fixed;
the upper plate moves with velocity \(Ue_x\). Seek a steady parallel flow
\[
  u=(V(y),0,0),\qquad p=\text{constant}.
\]
The incompressibility condition is automatic. The \(x\)-momentum equation is
\[
  0=\mu V''(y),
\]
with boundary conditions \(V(0)=0\), \(V(h)=U\). Thus
\[
  \boxed{V(y)=U\frac{y}{h}.}
\]
The shear stress is constant:
\[
  \sigma_{xy}=\mu V'(y)=\mu\frac{U}{h}.
\]

\subsection{Plane Poiseuille Flow}
\label{subsec:plane-poiseuille}

Now both plates are fixed and a pressure gradient drives the flow. Let
\[
  u=(V(y),0,0),
  \qquad
  -\frac{\dd p}{\dd x}=G>0,
  \qquad
  0<y<h.
\]
The steady \(x\)-momentum equation is
\[
  0=-\frac{\dd p}{\dd x}+\mu V''(y)=G+\mu V''(y),
\]
so
\[
  V''(y)=-\frac{G}{\mu}.
\]
With no-slip conditions \(V(0)=V(h)=0\),
\[
  \boxed{
  V(y)=\frac{G}{2\mu}y(h-y).
  }
\]
The maximum speed occurs at \(y=h/2\):
\[
  V_{\max}=\frac{Gh^2}{8\mu}.
\]
The volume flux per unit span is
\[
  Q=\int_0^h V(y)\,\dd y
  =
  \frac{Gh^3}{12\mu}.
\]

\subsection{Hagen-Poiseuille Pipe Flow}
\label{subsec:hagen-poiseuille}

For a circular pipe of radius \(a\), use cylindrical radius \(r\) and seek
\[
  u=V(r)e_z,\qquad -\frac{\dd p}{\dd z}=G>0.
\]
The axial equation is
\[
  0=G+\mu\frac1r\frac{\dd}{\dd r}
  \left(
    r\frac{\dd V}{\dd r}
  \right).
\]
Regularity at \(r=0\) and no slip at \(r=a\) give
\[
  \boxed{
  V(r)=\frac{G}{4\mu}(a^2-r^2).
  }
\]
The flux is
\[
  Q=2\pi\int_0^a V(r)r\,\dd r
  =
  \boxed{\frac{\pi a^4G}{8\mu}}.
\]
The fourth power in \(a\) is one of the most important practical consequences
of viscous laminar flow.

\begin{figure}[htbp]
\centering
\begin{tikzpicture}[scale=1.0, >=Latex]
  \draw[->] (-0.25,0) -- (-0.25,3.25) node[above] {\(y\)};
  \draw[->] (-0.25,0) -- (1.2,0) node[right] {\(x\)};
  \draw[thick] (0,0) -- (3.0,0);
  \draw[thick] (0,3.0) -- (3.0,3.0);
  \draw[->, thick, black!70] (2.1,3.18) -- (2.8,3.18) node[right] {wall speed \(U\)};
  \foreach \y/\len in {0.45/0.25,0.95/0.55,1.45/0.85,1.95/1.15,2.45/1.45}
    \draw[blue, thick, ->] (0.35,\y) -- ++(\len,0);
  \node[blue] at (1.5,3.45) {Couette};
  \node at (1.5,-0.35) {\(V(0)=0,\ V(h)=U\)};

  \begin{scope}[xshift=4.4cm]
    \draw[thick] (0,0) -- (3.0,0);
    \draw[thick] (0,3.0) -- (3.0,3.0);
    \draw[->, thick, black!70] (0.25,3.25) -- (1.35,3.25)
      node[right] {\(-\partial_xp=G\)};
    \foreach \y/\len in {0.45/0.45,0.95/1.05,1.50/1.45,2.05/1.05,2.55/0.45}
      \draw[red!75!black, thick, ->] (0.35,\y) -- ++(\len,0);
    \node[red!75!black] at (1.5,3.45) {Poiseuille};
    \node at (1.5,-0.35) {\(V(0)=V(h)=0\)};
  \end{scope}
\end{tikzpicture}
\caption{Couette flow is linear because it is driven by wall motion. Poiseuille
flow is parabolic because it balances pressure forcing against viscous
diffusion of momentum. The arrows are spatial velocity vectors, not a graph
with velocity on the vertical axis.}
\label{fig:couette-poiseuille-profiles}
\end{figure}

\section{Linear Stability Of Laminar Shear Flow}
\label{sec:linear-stability-shear}

Let a steady base flow be
\[
  U(y)e_x.
\]
Write a small perturbation as
\[
  u=U(y)e_x+\tilde u,\qquad
  p=P(x)+\tilde p,
\]
and retain only terms linear in \(\tilde u\). In nondimensional variables,
\[
  \boxed{
  \partial_t\tilde u
  +
  U\partial_x\tilde u
  +
  \tilde v\,U'(y)e_x
  =
  -\nabla\tilde p
  +
  \frac1{\mathrm{Re}}\Delta\tilde u,
  \qquad
  \nabla\cdot\tilde u=0.
  }
\]
The term \(\tilde v\,U'e_x\) is the mechanism by which cross-stream velocity
tilts fluid into faster or slower parts of the mean shear and extracts energy
from the base flow.

\begin{figure}[htbp]
\centering
\begin{tikzpicture}[scale=1.0, >=Latex]
  \draw[->] (-0.25,0) -- (-0.25,3.4) node[above] {\(y\)};
  \draw[->] (-0.25,0) -- (5.9,0) node[right] {\(x\)};
  \draw[thick] (0,0.35) -- (5.3,0.35);
  \draw[thick] (0,3.05) -- (5.3,3.05);
  \foreach \y/\len in {0.65/0.35,1.10/0.62,1.55/0.92,2.00/1.25,2.45/1.55,2.85/1.8}
    \draw[blue!70!black, thick, ->] (0.45,\y) -- ++(\len,0);
  \draw[red!75!black, thick, ->] (2.75,1.15) -- (2.75,2.05)
    node[midway,right] {\(\tilde v\)};
  \draw[red!75!black, thick, ->] (2.75,2.05) -- (3.55,2.05)
    node[right] {\(\tilde v\,U'(y)e_x\)};
  \draw[purple!75!black, thick, domain=0.3:5.0, samples=120]
    plot (\x,{1.75+0.22*sin(115*\x)});
  \node[blue!70!black] at (4.25,2.75) {base shear \(U(y)e_x\)};
  \node[purple!75!black] at (4.0,1.25) {normal-mode disturbance};
\end{tikzpicture}
\caption{Linear shear-flow stability. A wall-normal perturbation velocity
\(\tilde v\) moves fluid across the base shear; the term
\(\tilde v\,U'(y)e_x\) is the linear energy-exchange channel between the base
flow and the disturbance.}
\label{fig:shear-stability-mechanism}
\end{figure}

\subsection{Orr-Sommerfeld Equation}
\label{subsec:orr-sommerfeld}

For two-dimensional perturbations, introduce a streamfunction
\[
  \tilde u=\partial_y\psi,\qquad
  \tilde v=-\partial_x\psi.
\]
Seek normal modes
\[
  \psi(x,y,t)=\phi(y)e^{i\alpha x+\sigma t}.
\]
Taking the curl of the linearized equation removes pressure and gives the
linear vorticity equation. Since the perturbation vorticity is
\[
  \tilde\omega_z
  =
  \partial_x\tilde v-\partial_y\tilde u
  =
  -(\partial_y^2-\alpha^2)\phi\,e^{i\alpha x+\sigma t},
\]
substitution gives a fourth-order eigenvalue equation for the wall-normal
shape \(\phi(y)\):
\[
  \boxed{
  (\sigma+i\alpha U)(D^2-\alpha^2)\phi
  -
  i\alpha U''\phi
  =
  \frac1{\mathrm{Re}}(D^2-\alpha^2)^2\phi,
  }
\]
where \(D=\dd/\dd y\). No-slip and no-penetration at walls give
\[
  \phi=0,\qquad D\phi=0
\]
at the walls. This fourth-order eigenvalue problem is the
Orr-Sommerfeld equation. A base flow is linearly unstable if some eigenvalue has
\[
  \operatorname{Re}\sigma>0.
\]
Many fluid-mechanics texts write the same normal mode as
\(e^{i\alpha(x-ct)}\). With the convention used here, the complex phase speed
\(c\) and growth rate \(\sigma\) are related by
\[
  \sigma=-i\alpha c,
\]
so \(\operatorname{Re}\sigma=\alpha\,\operatorname{Im}c\) for real
\(\alpha\).

\begin{remark}[Classical facts]
Plane Poiseuille flow has a linear instability above a critical Reynolds number
about \(5772\), using centerline speed and half-channel width. Plane Couette
flow is linearly stable for all Reynolds numbers in the classical
Orr-Sommerfeld analysis, yet it can transition through finite-amplitude
disturbances. Stability is therefore not a single concept: linear eigenvalue
stability, energy stability, transient growth, and nonlinear transition answer
different questions.
\end{remark}

\subsection{Energy Stability Estimate}
\label{subsec:energy-stability}

For perturbations of a parallel flow, the kinetic energy of the disturbance
satisfies
\[
  \frac12\frac{\dd}{\dd t}\int_\Omega |\tilde u|^2\,\dd^3x
  =
  -\int_\Omega \tilde u\,\tilde v\,U'(y)\,\dd^3x
  -
  \frac1{\mathrm{Re}}
  \int_\Omega |\nabla\tilde u|^2\,\dd^3x.
\]
The second term is viscous dissipation. The first is production by mean shear.
If dissipation dominates production for all admissible perturbations, the base
flow is energy stable. If not, the estimate does not prove instability, but it
reveals the only possible energy source for perturbation growth.

\section{Point Vortices As A Hamiltonian Fluid Model}
\label{sec:point-vortices}

In two dimensions, an idealized vorticity distribution can be written as
\[
  \omega(x,y)=\sum_{i=1}^{N}\Gamma_i\delta(x-x_i,y-y_i),
\]
where \(\Gamma_i\) is circulation. The point vortex equations are
\[
\begin{aligned}
  \Gamma_i\dot x_i &= \frac{\partial H}{\partial y_i},\\
  \Gamma_i\dot y_i &= -\frac{\partial H}{\partial x_i},
\end{aligned}
\]
with Hamiltonian
\[
  H
  =
  -\frac{1}{4\pi}
  \sum_{i<j}\Gamma_i\Gamma_j
  \log\left((x_i-x_j)^2+(y_i-y_j)^2\right).
\]
This is Hamiltonian mechanics, but the symplectic form is weighted:
\[
  \omega_{\mathrm{vortex}}
  =
  \sum_i \Gamma_i\,\dd x_i\wedge\dd y_i.
\]
The weight \(\Gamma_i\) is the circulation carried by the \(i\)-th vortex. One
can derive this finite-dimensional symplectic form by reducing the ideal-fluid
Hamiltonian bracket to delta-function vorticity; here we take the reduced
point-vortex bracket as the model.

\begin{figure}[htbp]
\centering
\begin{tikzpicture}[scale=1.05, >=Latex]
  \draw[->] (-2.5,0) -- (2.7,0) node[right] {\(x\)};
  \draw[->] (0,-2.2) -- (0,2.4) node[above] {\(y\)};
  \foreach \x/\y/\g/\lab in {-0.65/-0.35/+/\Gamma,0.65/0.35/-/-\Gamma}
  {
    \draw[thick] (\x,\y) circle (0.16);
    \node at (\x,\y) {\(\g\)};
    \node[above right] at (\x,\y) {\(\lab\)};
  }
  \draw[blue, thick, ->] (-1.7,-0.35) -- (1.45,-0.35);
  \draw[blue, thick, ->] (-0.4,0.35) -- (2.75,0.35);
  \draw[dashed] (-0.65,-0.35) -- (0.65,0.35);
  \draw[<->] (-0.65,-0.62) -- (0.65,0.08) node[midway, below right] {separation};
  \node at (0,-1.75) {a vortex dipole translates by mutual induction};
\end{tikzpicture}
\caption{A vortex dipole is the simplest point-vortex motion: the two vortices
translate together with fixed separation. More complicated point-vortex
systems give finite-dimensional Hamiltonian models of ideal-fluid motion.}
\label{fig:vortex-dipole}
\end{figure}

\section{What To Take Forward}
\label{sec:fluid-take-forward}

\begin{itemize}
\item The material derivative is the bridge between following particles and
describing fields at fixed spatial points.
\item Conservation of mass and momentum are balance laws before they are
differential equations; the local PDEs come from applying the divergence theorem
to arbitrary moving volumes.
\item Navier-Stokes is Cauchy momentum balance plus the Newtonian stress law.
Changing the constitutive law changes the continuum model.
\item The Reynolds number compares inertial transport with viscous diffusion of
momentum. It is the organizing scale behind laminar solutions, instabilities,
and transition.
\item Point vortices show that Hamiltonian structure survives inside fluid
mechanics after a strong but explicit reduction.
\end{itemize}

\section{Associated Demos}
\label{sec:fluid-associated-demos}

\begin{itemize}
\item \path{demos/python/fluids_vorticity.py}, a point-vortex Hamiltonian model.
\item \path{demos/python/hamiltonian_pendulum.py}, useful for comparing
separatrix geometry with shear-flow stability.
\item Suggested extension: a numerical Orr-Sommerfeld spectrum for Couette and
Poiseuille flow.
\end{itemize}
